% BHU PYQ -- Manually transcribed & error-corrected
\documentclass[11pt,a4paper]{article}

\usepackage[a4paper,top=2.5cm,bottom=2.5cm,left=2.8cm,right=2.8cm,headheight=14pt]{geometry}
\usepackage{amsmath,amssymb}
\usepackage[T1]{fontenc}
\usepackage[utf8]{inputenc}
\usepackage{lmodern}
\usepackage{microtype}
\usepackage{fancyhdr}
\usepackage{enumitem}
\usepackage{hyperref}
\usepackage{lastpage}
\usepackage{parskip}

\hypersetup{colorlinks=true,urlcolor=black,linkcolor=black,
  pdftitle={CHB-601: Analytical Chemistry-I -- BHU PYQ},pdfauthor={Banaras Hindu University}}

\pagestyle{fancy}\fancyhf{}
\renewcommand{\headrulewidth}{0.4pt}\renewcommand{\footrulewidth}{0.4pt}
\fancyhead[L]{\small\textit{Banaras Hindu University}}
\fancyhead[R]{\small\textit{CHB-601: Analytical Chemistry-I}}
\fancyfoot[C]{\small Page \thepage\ of \pageref{LastPage}}

\newlist{parts}{enumerate}{2}
\setlist[parts,1]{label=(\alph*),leftmargin=*,itemsep=5pt,topsep=3pt}
\setlist[parts,2]{label=(\roman*),leftmargin=*,itemsep=3pt,topsep=2pt}
\newcommand{\pts}[1]{\hfill\textit{[#1]}}

\begin{document}

\begin{center}
  {\large\bfseries BANARAS HINDU UNIVERSITY}\\[2pt]
  {\normalsize B.Sc. (Hons.) Semester VI Examination, 2013-14}\\[6pt]
  \rule{0.8\linewidth}{0.5pt}\\[6pt]
  {\large\bfseries Paper No. CHB-601: Analytical Chemistry-I}\\[4pt]
  {\normalsize Time: 3 Hours\hfill Maximum Marks: 70}
\end{center}
\rule{\linewidth}{0.4pt}
\small
\noindent\textit{Instructions: Answer five questions in total, including Question~1 which is compulsory.}
\normalsize
\bigskip

\noindent\textbf{Question 1.} Answer all parts of the following: \pts{5 $\times$ 2 = 10}
\begin{parts}
  \item What are the units of molar absorptivity ($\varepsilon$) and specific absorbance ($a$)?
  \item Which chromatographic technique is a two-dimensional chromatography? Give one example.
  \item Define the term retention time ($t_R$) in chromatography.
  \item Beer's law is obeyed in the case of an aqueous solution of a weak acid when the solution is made strongly acidic or basic. Explain why.
  \item The distribution constant ($K_D$) for iodine between an organic solvent and water is $85$. Calculate the concentration of $\text{I}_2$ remaining in water after extraction.
\end{parts}

\medskip\rule{\linewidth}{0.2pt}

\noindent\textbf{Question 2.}
\begin{parts}
  \item Show that for $n$ successive extractions with volume $V_x$ of organic solvent, the fraction of solute remaining in aqueous phase of volume $V_{aq}$ is:
    \[ q = \left( \frac{V_{aq}}{V_{aq} + D V_x} \right)^n \] \pts{8}
  \item Forty milliliters of $0.20\,\text{M}$ butyric acid in water is extracted with $20\,\text{mL}$ of ether. After the layers are separated, it was found that $1.0\,\text{mmol}$ of butyric acid remained in the aqueous layer. Calculate the distribution ratio ($D$). \pts{6}
\end{parts}

\medskip\rule{\linewidth}{0.2pt}

\noindent\textbf{Question 3.} Illustrate the following solvent extraction techniques: \pts{3 $\times$ 4.5 = 13.5}
\begin{parts}
  \item Multiple extractions vs. single extraction
  \item Synergic solvent extraction
  \item Ion-pair extraction of metal chelates
\end{parts}

\medskip\rule{\linewidth}{0.2pt}

\noindent\textbf{Question 4.}
\begin{parts}
  \item A compound X exhibits a molar absorptivity of $2.45 \times 10^4\,\text{L}\,\text{mol}^{-1}\,\text{cm}^{-1}$ at $450\,\text{nm}$. What concentration of X in a solution causes a $25\%$ decrease in the radiant power at $450\,\text{nm}$ in a $1.0\,\text{cm}$ cell? \pts{6}
  \item What is a Ringbom plot (or Ayres plot)? Discuss its advantages and limitations in spectrophotometric analysis. \pts{8}
\end{parts}

\medskip\rule{\linewidth}{0.2pt}

\noindent\textbf{Question 5.}
\begin{parts}
  \item Discuss photometric errors in spectrophotometry and how they vary with transmittance ($T$). Prove that error is minimum when $T \approx 0.368$. \pts{8}
  \item Illustrate the principles of one-component and two-component analysis using UV-Visible spectrophotometry. \pts{6}
\end{parts}

\vfill
\rule{\linewidth}{0.4pt}
\begin{center}\small
  \href{https://bhu-exam.vercel.app/}{Click here to visit our website}\\[4pt]
  \href{https://me-aryan.vercel.app/}{Click here to visit our contributor}\\[4pt]
  \href{https://quashan-me.vercel.app/}{Click here to visit our contributor}
\end{center}

\end{document}
